% ============================================================================
% Ferro, C. A. T. (2014). Fair scores for ensemble forecasts.
% Q. J. R. Meteorol. Soc. 140: 1917--1923. doi:10.1002/qj.2270
%
% FULL-TEXT TRANSCRIPTION (no official LaTeX source exists). Transcribed by
% Claude from paper.pdf. Equation, definition, theorem and table numbers match
% the published paper. Figures are described, not reproduced; reference list
% omitted (see paper.pdf). Blocks marked [TRANSCRIBER NOTE] are NOT in the
% paper. Verify any equation you rely on against paper.pdf.
%
% ORIENTATION: all scores here are NEGATIVELY oriented (lower = better).
% Notation: ensemble x = (x_1,...,x_m), m members; observation y.
% ============================================================================
\documentclass{article}
\usepackage{amsmath,amssymb,amsthm}
\newtheorem*{definition1}{Definition 1}
\newtheorem*{definition2}{Definition 2}
\newtheorem*{definition3}{Definition 3}
\newtheorem*{definition4}{Definition 4}
\newtheorem*{theorem1}{Theorem 1}
\newtheorem*{example1}{Example 1}
\newtheorem*{example2}{Example 2}
\newcommand{\bx}{\mathbf{x}}
\newcommand{\dd}{\mathrm{d}}
\begin{document}

\title{Fair scores for ensemble forecasts}
\author{C. A. T. Ferro}
\date{}
\maketitle

\begin{abstract}
The notion of fair scores for ensemble forecasts was introduced recently to
reward ensembles with members that behave as though they and the verifying
observation are sampled from the same distribution. In the case of forecasting
binary outcomes, a characterization is given of a general class of fair scores
for ensembles that are interpreted as random samples. This is also used to
construct classes of fair scores for ensembles that forecast multicategory and
continuous outcomes. The usual Brier, ranked probability and continuous ranked
probability scores for ensemble forecasts are shown to be unfair, while adjusted
versions of these scores are shown to be fair. A definition of fairness is also
proposed for ensembles with members that are interpreted as being dependent and
it is shown that fair scores exist only for some forms of dependence.
\end{abstract}

\section{Introduction}
There are many ways to evaluate the performance of ensemble forecasts (e.g.\
Weigel, 2012). One approach converts ensembles to probability distributions and
then applies techniques for evaluating probability forecasts. This approach
evaluates not only the original ensemble, however, but also the
post-processing scheme used to create the probability forecast. We focus on
techniques for evaluating the ensemble alone. Existing techniques of this type
include rank histograms, conditional exceedance probabilities, generalized
discrimination and assessments of spread--skill relationships. As Weigel (2012)
notes, however, none of these techniques awards each individual ensemble its
own score; they measure attributes such as reliability that can be calculated
only for a set of (usually many) ensembles.

Scoring rules quantify the performance of each forecast individually. We follow
Murphy (1997) and define a scoring rule to be any function of a single forecast
and verifying observation, whatever the type of forecast. If $x$ is a scalar
point forecast and $y$ a verifying observation, a common scoring rule is the
squared error $(x-y)^2$. If $y$ is binary ($y=1$ when an event occurs, $y=0$
otherwise) and $p\in[0,1]$ is a probability forecast, then $(p-y)^2$ defines
the quadratic or Brier score (Brier, 1950). Both are \emph{negatively oriented}:
lower values correspond to better forecasts. Performance of a set of forecasts
is often summarized by the mean score.

A common idea is to apply scoring rules for probability forecasts to the
empirical distribution function of the ensemble members; e.g.\ the Brier score
with $p$ set to the proportion of members predicting the event. Regrettably,
this typically favours ensembles with members that behave as though they and
the verifying observation are sampled from \emph{different} distributions.
Example: let $\Pr(y=1)=1/4$ and let the ensemble have one binary member $x$.
Compare (a) $x$ an independent draw from the same distribution as $y$, and
(b) $x$ with $\Pr(x=1)=0$ (always zero; overconfident). The Brier score favours
(b): $E\{(0-y)^2\}=1/4$, whereas (a) receives $E\{(x-y)^2\}=3/8$.

Fricker \emph{et al.}\ (2013) introduced fair scoring rules to overcome this
problem. Section 2 formalizes fair scoring rules for ensembles interpreted as
random samples, characterises a wide class for binary outcomes, and uses it to
generate fair scores for multiple categories and continuous outcomes. Section 3
discusses ensembles with dependent members. Section 4 summarises.

\section{Independent ensemble members}
\subsection{Fair scoring rules}
A scoring rule $s(p,y)$ for probability forecasts is \emph{proper} if its
expectation with respect to any distribution $q$ for $y$ is optimized when
$p=q$. A scoring rule $s(x,y)$ for point forecasts is \emph{consistent} for a
functional $f$ if its expectation with respect to any $q$ is optimized when
$x=f(q)$ (e.g.\ squared error is consistent for the mean).

[Discussion, summarised:] If an ensemble is interpreted as the forecaster's
belief distribution (equal mass on each member) a proper score would be fair;
if interpreted as specific quantiles, a consistent score would be fair (e.g.\
Br\"ocker (2012): the CRPS of the empirical distribution function of an
$m$-member ensemble is consistent for the quantiles $(i-1/2)/m$,
$i=1,\dots,m$). More often, an ensemble is interpreted as a sample from a
distribution of possible outcomes. We interpret the ensemble $\bx$ as a random
sample whose members are independent realizations from some underlying
ensemble distribution. There are no scoring rules that \emph{elicit} random
samples, because the expectation of $s(\bx,y)$ w.r.t.\ $q$ is a deterministic
function of $\bx$, so the optimizing member values can always be issued instead.
Fricker \emph{et al.}\ (2013) proposed instead:

\begin{definition1}
A scoring rule, $s(\bx,y)$, is \emph{fair} for random sample ensembles, $\bx$,
if its expectation with respect to both the ensemble distribution, $p$, and any
distribution, $q$, for the verifying observation, $y$, is optimized when $p=q$.
The scoring rule is \emph{strictly fair} if its expectation is uniquely
optimized when $p=q$.
\end{definition1}

Equivalently: the expectation of $s(\bx,y)$ with respect to the ensemble
distribution $p$ is a (strictly) proper scoring rule for the probability
forecast $p$. Fair scores evaluate the underlying ensemble distribution; unfair
scores favour ensembles generated from imperfect distributions $p\neq q$.

The only fair scoring rules for ensembles with one member ($m=1$) are trivial:
given $s(x,y)$ and any $q$, let $X$ be the set of $x$ values optimizing the
expected score. Choosing $p$ with support contained in $X$ optimizes the
expectation; if $\mathrm{supp}(q)\not\subseteq X$ the expected score is not
optimized at $p=q$; if $\mathrm{supp}(q)\subseteq X$, any $p$ with the same
support as $q$ optimizes it.

Two restrictions:
\begin{definition2}
A scoring rule, $s(\bx,y)$, is \emph{ensemble-symmetric} if
$s(\bx,y)=s(\pi(\bx),y)$ for all ensembles $\bx$, verifying observations $y$,
and permutations $\pi(\bx)$ of $\bx$.
\end{definition2}
Equivalently, $s$ depends on $\bx=(x_1,\dots,x_m)$ only through its order
statistics, or only through the empirical distribution function
$\frac1m\sum_{i=1}^m I(x_i\le x)$, where $I(A)=1$ if $A$ is true and $0$
otherwise.

\begin{definition3}
A scoring rule, $s(\bx,y)$, is \emph{finite} if $|s(\bx,y)|<\infty$ for all
ensembles $\bx$ and verifying observations $y$.
\end{definition3}
This ensures the expectation w.r.t.\ the ensemble distribution is regular
(Gneiting and Raftery, 2007).

\subsection{Binary outcomes}
Let $x_1,\dots,x_m$ and $y$ take values in $\{0,1\}$; $p=\Pr(x_i=1)$,
$q=\Pr(y=1)$. Ensemble-symmetric rules depend on $\bx$ only through the number
of members equal to 1. Write $s_{i,j}$ for the score when $\sum_{l=1}^m x_l=i$
and $y=j$, $i=0,\dots,m$, $j=0,1$ (dependence on $m$ suppressed).

\begin{theorem1}
The negatively oriented, finite, ensemble-symmetric scoring rule defined by
$\{s_{i,j}: i=0,1,\dots,m;\ j=0,1\}$ is fair for random sample ensembles of
size $m$ if
\begin{equation}\tag{1}
  (m-i)(s_{i+1,0}-s_{i,0}) = i\,(s_{i-1,1}-s_{i,1})
\end{equation}
for $i=0,1,\dots,m$ and $s_{i+1,0}\ge s_{i,0}$ for $i=0,1,\dots,m-1$.
\end{theorem1}

Notes:
\begin{enumerate}
\item The undefined $s_{-1,1}$ ($i=0$) and $s_{m+1,0}$ ($i=m$) in (1) can be
ignored: read the RHS as zero when $i=0$ and the LHS as zero when $i=m$.
\item Theorem 1 remains true with ``positively oriented'' and
$s_{i+1,0}\le s_{i,0}$.
\item The equality constraints (1) are \emph{necessary} for fairness; the
inequalities $s_{i+1,0}\ge s_{i,0}$ (which via (1) imply
$s_{i+1,1}\le s_{i,1}$) are merely sufficient, but desirable.
\item If at least one inequality is strict, the expectation w.r.t.\ the
ensemble distribution is strictly proper for $p$, so $s(\bx,y)$ is strictly
fair. The only fair-but-not-strictly-fair rules in Theorem 1 are the trivial
ones with $s_{i,0}=s_{0,0}$ and $s_{i,1}=s_{m,1}$ for all $i$.
\item Fair rules are trivial when $m=1$: Theorem 1 requires $s_{1,0}=s_{0,0}$
and $s_{0,1}=s_{1,1}$.
\item Fairness requires $s_{1,0}=s_{0,0}$ and $s_{m-1,1}=s_{m,1}$: ensembles in
which all but one member is correct score the same as perfect forecasts
(necessary because of the discreteness of $\sum x_l$; otherwise a forecaster
could hedge by setting $p=0$ when $q$ is small).
\item The equality constraints (1) are the discrete analogue of the relationship
(Savage, 1971) satisfied by proper scoring rules $s(p,y)$ for continuous $p$:
\[
  (1-p)\frac{\dd}{\dd p}s(p,0) = -p\frac{\dd}{\dd p}s(p,1).
\]
Approximating $p$ by $i/m$, $\dd s(p,0)/\dd p$ by $m(s_{i+1,0}-s_{i,0})$ and
$\dd s(p,1)/\dd p$ by $m(s_{i,1}-s_{i-1,1})$ recovers (1).
\end{enumerate}

There are $m+1$ degrees of freedom in the fair scoring rules of Theorem 1.
Imposing complement symmetry ($s_{i,0}=s_{m-i,1}$), $s_{0,0}=s_{m,1}=0$ for
perfect forecasts and $s_{m,0}=s_{0,1}=1$ for completely wrong forecasts leaves
$\lfloor m/2\rfloor-1$ degrees of freedom when $m>1$.

\begin{table}[h]
\centering
\caption*{Table 1. Scores $s_{i,j}$ for negatively oriented fair scoring rules
when $m=2,3,4$. The constant $a$ may be any number in $[0,1/4]$.}
\begin{tabular}{c|cc|cc|cc}
 & \multicolumn{2}{c|}{$m=2$} & \multicolumn{2}{c|}{$m=3$} & \multicolumn{2}{c}{$m=4$}\\
$i$ & $j=0$ & $j=1$ & $j=0$ & $j=1$ & $j=0$ & $j=1$\\\hline
0 & 0 & 1 & 0 & 1 & 0 & 1\\
1 & 0 & 0 & 0 & 1/3 & 0 & $1-3a$\\
2 & 1 & 0 & 1/3 & 0 & $a$ & $a$\\
3 &   &   & 1 & 0 & $1-3a$ & 0\\
4 &   &   &   &   & 1 & 0
\end{tabular}
\end{table}

The usual ensemble Brier score (probability = proportion of members forecasting
the event) is
\begin{equation}\tag{2}
  s_{i,j} = \left(\frac{i}{m}-j\right)^2 .
\end{equation}
It does not satisfy (1), so it is \textbf{not fair}. Ferro (2007) and Ferro
\emph{et al.}\ (2008) proposed an adjusted Brier score, an unbiased estimator
of the score that would be obtained as ensemble size $\to\infty$:
\begin{equation}\tag{3}
  s_{i,j} = \left(\frac{i}{m}-j\right)^2 - \frac{i(m-i)}{m^2(m-1)} ,
\end{equation}
which satisfies Theorem 1 and so \textbf{is fair} for $m>1$. It is
complement-symmetric, scores 0 and 1 for perfect and completely wrong
forecasts, lies in $[0,1]$, and is $\le$ (2) with equality iff $i=0$ or $m$:
the original score unfairly penalises ensembles by not accounting for finite
ensemble size.

\emph{Figure 1} [described]: expected adjusted (black) vs original (grey)
Brier scores for $m=2,4,8$ against $p$, for $q=1/2$ and $q=1/4$. The expected
original score is minimised at $p\neq q$ unless $q\in\{0,1/2,1\}$; the expected
adjusted score is always minimised at $p=q$ (and does not depend on $m$).
The original score favours overconfident ensembles: for $q=0.25$ its
expectation is optimised at $p=0$ ($m=2$), $p=0.17$ ($m=4$), $p=0.21$ ($m=8$).
In a seasonal precipitation example (Ferro 2007) adjusted scores are typically
$\sim$5\% smaller; rankings rarely change.

\subsection{Multiple categories and continuous outcomes}
\textbf{Multicategory.} Suppose $n$ possible outcomes $1,\dots,n$; ensemble
distribution $(p_1,\dots,p_n)$ with $p_k=\Pr(x_i=k)$, observation distribution
$(q_1,\dots,q_n)$ with $q_k=\Pr(y=k)$. Let $y_k=I(y=k)$ and
$\bx_k=(x_{k,1},\dots,x_{k,m})$ with $x_{k,i}=I(x_i=k)$. For each $k$ let
$s_k(\bx_k,y_k)$ be a fair, finite, ensemble-symmetric, negatively oriented
scoring rule as in Theorem 1. Then
\begin{equation}\tag{4}
  s(\bx,y) = \sum_{k=1}^n s_k(\bx_k,y_k)
\end{equation}
is a fair, finite, ensemble-symmetric scoring rule for $\bx$. If
$s_k(\bx_k,y_k)=s^k_{i,j}$ when $\sum_l x_{k,l}=i$, $y_k=j$, we need only that
$s^k_{i,j}$ satisfy Theorem 1 for each $k$. Taking $s^k_{i,j}$ = adjusted Brier
score (3) gives the adjusted multicategory Brier score of Ferro \emph{et al.}\
(2008); redefining $y_k=I(y\le k)$, $x_{k,i}=I(x_i\le k)$ gives their fair
adjusted ranked probability score.

\textbf{Continuous outcomes.} Let ensemble members and observation be real;
$p$ and $q$ the densities of the ensemble distribution and of $y$. Let
$y_t=I(y\le t)$ and $\bx_t=(x_{t,1},\dots,x_{t,m})$, $x_{t,i}=I(x_i\le t)$:
the binary observation and ensemble forecast for the event $\{y\le t\}$. For
each threshold $t$ let $s_t(\bx_t,y_t)$ be fair, finite, ensemble-symmetric,
negatively oriented (Theorem 1). If $|s_t(\bx_t,y_t)|$ is bounded for all
$\bx_t,y_t,t$ and $r$ is any probability density on $\mathbb{R}$, then by
Fubini's Theorem
\begin{equation}\tag{5}
  s(\bx,y) = \int_{-\infty}^{\infty} s_t(\bx_t,y_t)\,r(t)\,\dd t
\end{equation}
is a fair, finite, ensemble-symmetric scoring rule for continuous outcomes.
Taking $s_t$ equal to the adjusted Brier score (3) (with $\sum_l x_{t,l}=i$,
$y_t=j$) makes (5) a weighted version ($r$ weighting thresholds) of the
\textbf{adjusted (fair) CRPS} of Ferro \emph{et al.}\ (2008) and Fricker
\emph{et al.}\ (2013), which may be written
\[
  s(\bx,y) = \sum_{i=1}^{k+1} s_{i-1,0}\int_{z_{(i)}}^{z_{(i+1)}} r(t)\,\dd t
  \;+\; \sum_{i=k+1}^{m+1} s_{i-1,1}\int_{z_{(i+1)}}^{z_{(i+2)}} r(t)\,\dd t ,
\]
where $z_{(i)}$ is the $i$th order statistic of the set
$\{-\infty,x_1,\dots,x_m,y,\infty\}$ and $k$ is such that $z_{(k+2)}=y$. The
unweighted version of Ferro \emph{et al.}\ (2008) is $r(t)=1$ for all $t$.
Substituting the original Brier score (2) for $s_t$ gives the original
(unfair) ensemble CRPS.

\paragraph{[TRANSCRIBER NOTE --- not in the paper] Kernel form.}
With $r\equiv1$ and $F_m(t)=\frac1m\sum_i I(x_i\le t)$, (3) gives
$s_t=(F_m(t)-y_t)^2 - F_m(t)(1-F_m(t))/(m-1)$. Using
$\int(F_m-I(y\le t))^2\dd t=\frac1m\sum_i|x_i-y|-\frac1{2m^2}\sum_{i,j}|x_i-x_j|$
and $\int F_m(1-F_m)\,\dd t=\frac{1}{2m^2}\sum_{i,j}|x_i-x_j|$:
\[
  \mathrm{CRPS}(\bx,y)=\frac1m\sum_{i=1}^m|x_i-y|-\frac{1}{2m^2}\sum_{i,j=1}^m|x_i-x_j|,
  \qquad
  \mathrm{fCRPS}(\bx,y)=\frac1m\sum_{i=1}^m|x_i-y|-\frac{1}{2m(m-1)}\sum_{i,j=1}^m|x_i-x_j|,
\]
matching AIFS-CRPS eqs (1)--(2). $E_p[\mathrm{fCRPS}]=E|X-y|-\tfrac12E|X-X'|$
(the population CRPS of $p$), since $\frac{1}{m(m-1)}\sum_{i\ne j}|x_i-x_j|$ is
unbiased for $E|X-X'|$; the standard ensemble CRPS instead has expectation
$E|X-y|-\frac{m-1}{2m}E|X-X'|$, rewarding under-dispersion.

\emph{Figure 2} [described]: expected adjusted (black) vs original (grey)
CRPS for $m=2,4,8$ when $y\sim N(0,\beta^2)$ and $x_i$ are i.i.d.\
$N(0,\alpha^2)$, plotted against $\alpha$ for $\beta=1$ (solid) and
$\beta=1/2$ (dashed). The expectation of the original CRPS (computable from
Gneiting and Raftery, 2007) is minimised at $\alpha<\beta$; the adjusted CRPS
is always minimised at $\alpha=\beta$ (and does not depend on $m$). The original
CRPS favours overconfident (under-dispersed) ensembles: for $\beta=1$ it is
minimised at $\alpha=0.38$ ($m=2$), $0.63$ ($m=4$), $0.79$ ($m=8$).

\begin{table}[h]
\centering
\caption*{Table 2. Original and adjusted CRPS (mm\,month$^{-1}$, estimated
standard errors in brackets) for two seasonal precipitation forecasts.}
\begin{tabular}{lcc}
 & Original & Adjusted\\\hline
ECMWF & 51 (8.8) & 49 (8.8)\\
M\'et\'eo-France & 42 (6.5) & 39 (6.5)
\end{tabular}
\end{table}
Adjusted scores are about 5\% smaller; ordering unchanged; differences small
relative to standard errors.

\section{Dependent ensemble members}
Interpreting ensembles as random samples is often a good approximation, but we
may be told, or choose, to interpret members as dependent. If the $m$ members
are not i.i.d., consider their joint distribution $p_m$. Since $q$ is
one-dimensional we cannot require $p_m=q$. We consider only
\emph{exchangeable} members (joint distribution symmetric in its $m$
arguments; includes independence), each with the same marginal $p$.

\begin{definition4}
A scoring rule, $s(\bx,y)$, is \emph{fair for exchangeable ensembles}, $\bx$,
if its expectation with respect to both the joint distribution of the ensemble
members, $p_m$, and any distribution, $q$, for the verifying observation, $y$,
is optimized when $p=q$, where $p$ is the marginal distribution of $p_m$.
\end{definition4}

For $m=1$ only trivial fair rules exist; perfectly dependent members behave as
one member. Even without perfect dependence, only trivial fair rules may exist:

\begin{example1}
Binary outcomes, $m=2$. Let $q=\Pr(y=1)$, $p=\Pr(x_1=1)=\Pr(x_2=1)$ the
marginal, and $\Pr(x_1=1,x_2=1)=p^c$ with constant $1<c<\infty$ ($c=2$:
independent). With ensemble-symmetric $s_{i,j}$ ($x_1+x_2=i$, $y=j$), fairness
requires the expected score
\[
  (1-q)\{p^c s_{2,0}+2(p-p^c)s_{1,0}+(1-2p+p^c)s_{0,0}\}
  + q\{p^c s_{2,1}+2(p-p^c)s_{1,1}+(1-2p+p^c)s_{0,1}\}
\]
to be optimized at $p=q$. If $c\neq2$ (dependent members), this holds for
$0<q<1$ only if $s_{0,0}=s_{1,0}=s_{2,0}$ and $s_{0,1}=s_{1,1}=s_{2,1}$:
i.e.\ only for trivial rules that do not depend on the ensemble.
\end{example1}

\begin{example2}
Binary outcomes, $q=\Pr(y=1)$, $p=\Pr(x_i=1)$, and
$\mathrm{corr}(x_i,x_j)=c<1$ for all $i\neq j$, with $c$ independent of $p$.
Then the negatively oriented scoring rule
\begin{equation}\tag{6}
  s_{i,j} = \left(\frac{i}{m}-j\right)^2 - \left(1+\frac{cm}{1-c}\right)\frac{i(m-i)}{m^2(m-1)}
\end{equation}
is fair. (Proof: its expected value is $p^2-2pq+q$, minimised at $p=q$.) When
$c=0$ it reduces to the adjusted Brier score (3), which is therefore fair for
ensembles with pairwise \emph{uncorrelated} members, not only independent ones.
\end{example2}

No single (non-trivial) scoring rule is fair for all exchangeable ensembles: we
must specify the dependence structure and check whether a fair rule exists. If
none exists, a ``nearly fair'' rule (expectation optimised at $p$ within a small
distance of $q$) could be used; left for future research. The dependence
structure is rarely stated, and a fair rule cannot be chosen from a dependence
structure estimated from the ensemble being verified (e.g.\ replacing $c$ in
(6) by an estimate breaks fairness). Recommendation: verify ensembles under one
or more interpretations of interest; if an ensemble will be used as a random
sample, use the fair scores of Section 2; if its empirical distribution will be
used as a probability forecast, use proper scores on the empirical distribution
(typically an inappropriate interpretation).

\section{Summary and discussion}
A scoring rule for an ensemble forecast is fair if the expectation of the score
with respect to the distributions of both the ensemble members and the
verifying observation is optimized when these distributions coincide. With one
member, the only fair rules are trivial, so verification measures may favour
under-dispersed ensembles. Fair rules exist for $m>1$ i.i.d.\ members, for which
ensemble-symmetric rules are argued for. For binary outcomes a general class of
fair rules is characterised (Theorem 1), including an adjusted ensemble Brier
score; classes of fair rules for multicategory and continuous outcomes are
constructed, including adjusted RPS and CRPS. Fair rules can exist for
dependent members but are specific to the dependence structure and do not exist
for some forms of dependence. Since the correct dependence structure is
typically unknown, ensembles should be verified using scores fair for
dependence structures of interest, including independence.

\appendix
\section*{Appendix: proof of Theorem 1}
We need conditions on $s_{i,j}$ making $s(\bx,y)$ fair, i.e.\ making its
expectation w.r.t.\ the ensemble distribution a proper scoring rule for $p$.
With independent members, $\Pr(x_i=1)=1-\Pr(x_i=0)=p$, the expectation is
\begin{equation}\tag{A1}
  s_m(p,y) = \sum_{i=0}^m \binom{m}{i} p^i(1-p)^{m-i} s_{i,y},
\end{equation}
regular since $|s_{i,y}|<\infty$. Following Savage (1971), Gneiting and Raftery
(2007) showed that any regular, negatively oriented scoring rule $s(p,y)$ is
proper iff
\begin{equation}\tag{A2}
  s(p,y) = G(p) + (y-p)G'(p)
\end{equation}
for a concave $G$ (strictly concave for strict propriety), $G'$ its derivative
if differentiable. With $\Pr(y=1)=q$,
\[
  s_m(p,q) = \sum_{i=0}^m\binom{m}{i}p^i(1-p)^{m-i}\{(1-q)s_{i,0}+q\,s_{i,1}\},
  \qquad s(p,q) = G(p)+(q-p)G'(p).
\]
Evaluating at $q=p$ requires
\begin{equation}\tag{A3}
  G(p) = \sum_{i=0}^m\binom{m}{i}p^i(1-p)^{m-i}\{(1-p)s_{i,0}+p\,s_{i,1}\},
\end{equation}
so $G$ is differentiable. For propriety, $s_m(p,q)$ must be minimised at $p=q$
for all $0\le q\le1$. Its derivative w.r.t.\ $p$ at $p=q$ is
\[
  \sum_{i=0}^m\binom{m}{i}q^i(1-q)^{m-i}\{(m-i)(s_{i+1,0}-s_{i,0}) - i(s_{i-1,1}-s_{i,1})\},
\]
which is zero for all $0<q<1$ iff the equality constraints (1) hold; these are
therefore necessary. Differentiating (A3),
\[
  G'(p) = \sum_{i=0}^m\binom{m}{i}\{(1-p)s_{i,0}+p\,s_{i,1}\}\frac{\dd}{\dd p}p^i(1-p)^{m-i}
  + \sum_{i=0}^m\binom{m}{i}p^i(1-p)^{m-i}(s_{i,1}-s_{i,0}).
\]
The first sum is the derivative of $s_m(p,q)$ w.r.t.\ $p$ at $q=p$, zero under
(1), leaving
\begin{equation}\tag{A4}
  G'(p) = \sum_{i=0}^m\binom{m}{i}p^i(1-p)^{m-i}(s_{i,1}-s_{i,0}).
\end{equation}
Hence $G(p)+(y-p)G'(p)=s_m(p,y)$, the form (A2); it remains to find conditions
for $G$ concave. By (A4), $G'(p)=s_m(p,1)-s_m(p,0)$. If $s_{i+1,0}\ge s_{i,0}$
for $i=0,\dots,m-1$, then (1) implies $s_{i+1,1}\le s_{i,1}$ as well; so as $p$
increases $s_m(p,1)$ decreases and $s_m(p,0)$ increases, so $G'(p)$ decreases,
i.e.\ $G$ is concave. If $s_{i+1,0}>s_{i,0}$ for at least one $i$, $G$ is
strictly concave. \qed

\end{document}
